%************************************************
\chapter{The Long Road}\label{app:long-road}
%************************************************

The ideas of this book stand at the end of a very long road. This appendix is a map of that road. It follows four threads through fifty-two \emph{stations}, places and moments where an idea changed, and thirteen \emph{meditations}, pauses to ask what the change meant. Its guiding thought is this:
\begin{quote}
\itshape In each of these histories a question was asked that the available concepts could not answer, and the age concluded that the question was defective. In each of them somebody eventually concluded instead that the concepts were defective, stipulated a new object, endured the ridicule, and handed the next generation a world with more in it.
\end{quote}

The tables below list the stations (numbered 1 to 52) and the meditations (numbered I to XIII) in order. Several of them appear as scenes in the chapters of this book.

%------------------------------------------------
\section*{Part I \quad Number}
%------------------------------------------------

How counting became arithmetic, and how each new kind of number, from fractions to the complex numbers, was forced on its inventors by a question the old numbers could not answer (Chapter~\ref{ch:one-number}).

\begin{center}
\small
\begin{tabularx}{\textwidth}{@{}r >{\raggedright\arraybackslash}p{0.36\textwidth} >{\raggedright\arraybackslash}X@{}}
\toprule
 & \tableheadline{Title} & \tableheadline{Place and time} \\
\midrule
1 & The Notch & southern Africa, roughly forty thousand years ago \\
I & The Limit and the Unlimited & Miletus and Croton, sixth and fifth centuries BCE \\
2 & The Word & everywhere, before writing \\
3 & The Clay & Sumer and Babylon, about 3300 to 1600 BCE \\
4 & The Rope and the Fraction & Egypt, about 1850 BCE \\
5 & The Crisis & Greece, fifth century BCE \\
6 & The Debt and the Void & China, about 200 BCE; India, 628 CE \\
7 & The House of Wisdom & Baghdad, ninth century \\
8 & The Long Ride West & B\'eja\"ia to Pisa, 1202 \\
9 & The Duel & Bologna and Milan, 1515 to 1545 \\
10 & The Wild Thought & Bologna, 1572 \\
II & The Freedom to Stipulate & from Euclid's definitions to the imaginary unit \\
11 & Two Centuries of Embarrassment & Europe, 1572 to 1777 \\
III & The Amphibian & Europe, 1700 to 1830 \\
12 & The Picture & Copenhagen 1797, Geneva 1806, G\"ottingen 1831 \\
13 & The Legitimation & Dublin, 1837 \\
14 & The End of the Road & and why it ends there \\
\bottomrule
\end{tabularx}
\end{center}

%------------------------------------------------
\section*{Part II \quad Angle}
%------------------------------------------------

How the shadow of a stick became trigonometry: the right angle of the builders' rope, the tables of chords of the Greek astronomers, the Indian half-chord that became our sine through a chain of translations, and the moments when an angle became a number and a direction became a quantity (Chapter~\ref{ch:oblique}).

\begin{center}
\small
\begin{tabularx}{\textwidth}{@{}r >{\raggedright\arraybackslash}p{0.36\textwidth} >{\raggedright\arraybackslash}X@{}}
\toprule
 & \tableheadline{Title} & \tableheadline{Place and time} \\
\midrule
15 & The Shadow & Egypt and Babylon, before 1500 BCE \\
16 & The Rope and the Right Angle & the Near East, second millennium BCE \\
17 & The Triangle Becomes an Argument & Ionia and Alexandria, sixth to third century BCE \\
18 & The Chord & Rhodes and Alexandria, 150 BCE to 150 CE \\
IV & Is an Angle a Quality, a Quantity, or a Relation? & Alexandria and Athens, second to fifth century CE \\
19 & The Half-Chord & India, fifth to fifteenth century \\
20 & A Word Loses Its Way & from Sanskrit to Latin \\
21 & The Discipline Comes of Age & Baghdad, Isfahan and Mar\=agha, ninth to thirteenth century \\
22 & The Angle Becomes a Number & Europe, 1637 to 1748 \\
23 & The Hinge & the sentence this part was written for \\
24 & The Direction Becomes a Quantity & Copenhagen 1797, Dublin 1843, New Haven 1901 \\
25 & The Payoff & Paris 1822, Chicago 1893 \\
\bottomrule
\end{tabularx}
\end{center}

%------------------------------------------------
\section*{Part III \quad Matter and Force}
%------------------------------------------------

How the stuff of the world became a quantity: from the atoms of the Greek and Indian philosophers and the balance of Archimedes, through Galileo's falling bodies and Newton's quantity of matter, to the two inverse square laws and the units that finally let both laws be written as equations (Chapters~\ref{ch:one-number} and~\ref{ch:interaction}).

\begin{center}
\small
\begin{tabularx}{\textwidth}{@{}r >{\raggedright\arraybackslash}p{0.36\textwidth} >{\raggedright\arraybackslash}X@{}}
\toprule
 & \tableheadline{Title} & \tableheadline{Place and time} \\
\midrule
26 & The Cut & Abdera and India, fifth century BCE \\
V & The Cut & Elea, Abdera and Athens, fifth and fourth centuries BCE \\
27 & The Rejection & Athens, fourth century BCE \\
VI & Abstraction & Athens, fourth century BCE \\
28 & The Balance & Syracuse, third century BCE \\
29 & The Push That Stays & Alexandria, Baghdad and Paris, sixth to fourteenth century \\
30 & The Atom Returns & Europe, 1417 to 1661 \\
31 & Fall & Padua and Arcetri, 1590 to 1638 \\
32 & Quantity of Matter & Cambridge, 1687 \\
33 & The First Inverse Square & London and Cambridge, 1679 to 1687 \\
VII & Necessity or Habit & Baghdad, C\'ordoba and Edinburgh, 1095 to 1748 \\
34 & Weighing the Earth & London, 1798 \\
35 & The Other Force & from amber to the Leyden jar, 600 BCE to 1750 CE \\
36 & The Second Inverse Square & Leeds, London and Paris, 1767 to 1785 \\
37 & The Units & Paris 1799, G\"ottingen 1832, Chicago 1893, Paris 1889 and 1948 \\
VIII & Force as a Fiction, Mass as a Ratio & Dublin 1721, Prague 1883 \\
38 & The Atom Becomes a Number & Manchester 1808, Paris 1913, Chicago 1913 \\
IX & From Substance to Function & Berlin, 1910 \\
39 & The Definitions Change Places & Versailles, 20 May 2019 \\
\bottomrule
\end{tabularx}
\end{center}

%------------------------------------------------
\section*{Part IV \quad The Language}
%------------------------------------------------

How mathematics came to be seen as the language of nature, from the Pythagoreans to Wigner, together with the doubts and the explanations that followed every such claim (Chapter~\ref{ch:prologue}).

\begin{center}
\small
\begin{tabularx}{\textwidth}{@{}r >{\raggedright\arraybackslash}p{0.36\textwidth} >{\raggedright\arraybackslash}X@{}}
\toprule
 & \tableheadline{Title} & \tableheadline{Place and time} \\
\midrule
40 & All Is Number & Croton, fifth century BCE \\
41 & God Ever Geometrizes & Athens, fourth century BCE \\
42 & The Door and the Key & Oxford and Paris, thirteenth century \\
43 & The Book of Nature & Rome, 1623 \\
44 & Geometry Shining in the Mind of God & Graz, Prague and Linz, 1596 to 1619 \\
45 & I Frame No Hypotheses & Cambridge, 1687 and 1713 \\
46 & As Much Science as There Is Mathematics & K\"onigsberg 1786, Paris 1822 \\
47 & Translate It Out of Hieroglyphics & London and Aberdeen, 1857 to 1865 \\
48 & Wiser Than Their Discoverers & Karlsruhe, 1887 \\
49 & Express It in Numbers & London, 1883 \\
X & Kant, Helmholtz, and the Enlargement of Perception & K\"onigsberg 1781, Heidelberg 1870 \\
50 & Certain, or About Reality & Berlin, 1921 \\
51 & Beauty in the Equations & Cambridge, 1928 to 1963 \\
XI & Labels on Representations & Princeton, 1939 \\
XII & Why They Have Not Merged & 1919 to 1975 \\
52 & Unreasonable Effectiveness & New York, 1960 \\
XIII & The Ladder & what each new number bought \\
\bottomrule
\end{tabularx}
\end{center}

